\documentclass{beamer}
\usetheme{Madrid}

% Encoding and fonts for pdfLaTeX
\usepackage[utf8]{inputenc}
\usepackage[T1]{fontenc}
\usepackage{lmodern}
\usepackage{hyperref}

\title{Prompt Engineering}
\author{Mehmet Ali Akyol, PhD}
\institute{Ankara Medipol University}
\date{\today}

\begin{document}

\begin{frame}
    \titlepage
    \begin{center}
        \vspace{1cm}
        \textbf{Week 9: Multimodal Prompting}
    \end{center}
\end{frame}

\section{Overview}

\begin{frame}
    \frametitle{Weekly Learning Objectives}
    \begin{itemize}
        \item Understand capabilities and limits of popular multimodal models.
        \item Craft prompts that fuse text with images, tables, or audio transcripts.
        \item Plan verification strategies for non-text outputs.
        \item Design accessibility-aware workflows (alt text, captions, transcripts).
        \item Evaluate ethical considerations in multimodal use.
    \end{itemize}
\end{frame}

\begin{frame}
    \frametitle{Introduction to Multimodality}
    \begin{block}{One-sentence intuition}
        	extbf{Modality} means the \textit{type of input}: text, image, audio, video, sensor data.
        	extbf{Multimodality} means the model can use \textit{more than one} at the same time.
    \end{block}
    \begin{itemize}
        \item \textbf{Defining Modality:} A specific mode of communication or sensory channel (e.g., text, vision, audio, sensor data).
        \item \textbf{Human Perception:} Humans are inherently multimodal. We don't just read; we see, hear, and feel to build a complete world model.
        \item \textbf{The Digital Evolution:}
        \begin{itemize}
            \item \textit{Web 1.0/2.0:} Primarily text-based (Search, Blogs).
            \item \textit{Modern Web:} Rich media (Video, Images, Voice).
        \end{itemize}
        \item \textbf{The Gap:} Until recently, AI treated these modalities separately (CV for images, NLP for text). Multimodality bridges them.
    \end{itemize}
\end{frame}

\begin{frame}
    \frametitle{Where do we actually use multimodal prompts?}
    \begin{itemize}
        \item \textbf{Photos:} A product photo, a broken device, a street scene.
        \item \textbf{Screenshots:} App UI bugs, error dialogs, settings pages.
        \item \textbf{Documents:} PDF pages, invoices, lab reports, forms.
        \item \textbf{Charts/Tables:} A plot in a paper, a dashboard screenshot.
        \item \textbf{Audio/Video:} Meeting recordings, lectures (often via transcript + key frames).
    \end{itemize}
    \begin{block}{Simple idea}
        Your prompt should say \textbf{what to look at} and \textbf{what to produce} (summary, table, JSON, steps, etc.).
    \end{block}
\end{frame}

\begin{frame}
    \frametitle{What is Multimodal AI?}
    \begin{block}{Definition}
        \textbf{Multimodal AI} refers to systems that can process, understand, and generate information using multiple types of data (modalities) simultaneously, such as text, images, audio, and video.
    \end{block}
    \begin{itemize}
        \item \textbf{Unimodal:} Text-in $\to$ Text-out (e.g., original ChatGPT).
        \item \textbf{Multimodal:} Image+Text-in $\to$ Text-out (e.g., GPT-4o, Gemini).
        \item \textbf{Goal:} Mimic human perception by integrating diverse sensory inputs into a unified latent space.
    \end{itemize}
\end{frame}

\begin{frame}
    \frametitle{Two common multimodal tasks}
    \begin{itemize}
        \item \textbf{Understanding:} "Describe what's in the image", "Extract fields from this invoice", "Explain the chart trend".
        \item \textbf{Decision support:} "Given this screenshot and error log, propose likely causes and next debugging steps".
        \item \textbf{Generation (less common in class prompts):} "Generate alt-text" or "Draft a report based on the image".
    \end{itemize}
    \begin{block}{}
        If you want \textbf{accuracy}, ask for \textbf{observations first}, then ask for \textbf{interpretation}.
    \end{block}
\end{frame}

\begin{frame}
    \frametitle{Why Multimodal?}
    \begin{itemize}
        \item \textbf{Information Density:} An image is worth 1000 tokens. Visuals capture spatial relationships that text struggles to describe.
        \item \textbf{Real-World Interaction:} The world is not just text. Robots, assistants, and tools need to ``see'' and ``hear.''
        \item \textbf{Emergent Capabilities:} Reasoning across modalities (e.g., explaining a meme) requires understanding both the visual content and the cultural text context.
    \end{itemize}
    \begin{block}{Mini examples}
        \begin{itemize}
            \item "Here is a screenshot of my error. Explain what this message means and what to try next."
            \item "Here is a graph. Is the trend increasing, decreasing, or seasonal?"
        \end{itemize}
    \end{block}
\end{frame}

\begin{frame}
    \frametitle{Key Concepts}
    \begin{itemize}
        \item \textbf{Vision-Language Models (VLMs):} Models trained on massive datasets of image-text pairs to learn associations.
        \item \textbf{Grounding:} The ability of the model to link textual concepts to specific pixels or regions in an image (e.g., drawing a bounding box around ``the cat'').
        \item \textbf{Alignment:} Ensuring the model's visual understanding matches human intent and safety standards.
        \item \textbf{Modality Gap:} The disconnect that can occur when a model is strong in text but weak in vision (or vice versa).
    \end{itemize}
\end{frame}

\begin{frame}
    \frametitle{Key Concepts (continued)}
    \begin{itemize}
        \item \textbf{OCR vs. Understanding:} Reading text in an image is OCR; explaining what it means is understanding.
        \item \textbf{Layout Understanding:} Tables, forms, headers/footers matter (document AI).
        \item \textbf{Evidence:} Good prompts ask the model to \textit{quote} the text it read or point to the region.
        \item \textbf{Uncertainty:} For hard images (tiny text), ask for "best guess" + "what is unclear".
    \end{itemize}
\end{frame}

\section{Capabilities and Limits}

\begin{frame}
    \frametitle{How VLMs Work (Simplified)}
    \begin{itemize}
        \item \textbf{Tokenization:} Just as text is broken into tokens, images are broken into ``patches'' (e.g., 16x16 pixel squares).
        \item \textbf{Encoding:} A Vision Encoder (like ViT) converts these patches into vectors.
        \item \textbf{Projection:} A projection layer maps these visual vectors into the same ``language space'' as text tokens.
        \item \textbf{Generation:} The LLM processes the sequence of [Image Vectors] + [Text Tokens] to predict the next text token.
    \end{itemize}
    \begin{block}{Analogy}
        Think of an image as being converted into "visual tokens" so the model can process it similarly to words.
    \end{block}
\end{frame}

\begin{frame}
    \frametitle{Model Landscape}
    \begin{itemize}
        \item \textbf{Frontier Models (Closed) --- examples (2025 era):}
        \begin{itemize}
            \item \textbf{OpenAI GPT family (e.g., GPT-4o/4.1-class):} Strong image understanding + natural language interaction.
            \item \textbf{Anthropic Claude 3.x:} Strong long-context reasoning; useful for document + image workflows.
            \item \textbf{Google Gemini 1.5/2.0-class:} Strong multimodal context handling and tool integrations.
        \end{itemize}
        \item \textbf{Open Weights / Source --- examples:}
        \begin{itemize}
            \item \textbf{Llama vision-capable releases:} Popular for building prototypes with local inference.
            \item \textbf{Qwen VL-family:} Strong document understanding and OCR-style tasks.
            \item \textbf{DeepSeek VL-family:} Good performance/cost trade-offs for vision+language tasks.
        \end{itemize}
    \end{itemize}
\end{frame}


\begin{frame}
    \frametitle{Failure Modes \& Risks}
    \begin{itemize}
        \item \textbf{Visual Hallucination:} Describing objects that aren't there. (e.g., ``The man is wearing a red hat'' when it's blue).
        \item \textbf{OCR Errors:} Misreading small text, handwriting, or rotated text in images.
        \item \textbf{Spatial Confusion:} Struggling with ``left of'', ``behind'', or counting overlapping objects.
        \item \textbf{Prompt Injection (Visual):} Hidden text in an image (steganography) that overrides system instructions (e.g., white text on white background saying ``Ignore safety rules'').
    \end{itemize}
    \begin{block}{Practical mitigations}
        \begin{itemize}
            \item Ask for \textbf{observations} first: "List what you can see".
            \item For OCR: ask it to \textbf{quote the text} it read.
            \item For counting: ask for a \textbf{count + uncertainty} ("I see about N").
        \end{itemize}
    \end{block}
\end{frame}

\section{Prompting Techniques}

\begin{frame}
    \frametitle{A simple multimodal prompt recipe (ROLE--CONTEXT--TASK--OUTPUT)}
    \begin{itemize}
        \item \textbf{Role:} Who should the model be? ("Act as a radiology assistant", "Act as a QA engineer")
        \item \textbf{Context:} What is the image/document? What should it focus on?
        \item \textbf{Task:} What exactly should it do? (extract, explain, compare, check)
        \item \textbf{Output:} Format + constraints (bullet list, table, JSON; "If uncertain, say so")
    \end{itemize}
    \begin{block}{Key trick}
        Add a line: \textit{"First describe what you observe. Then interpret."}
    \end{block}
\end{frame}

\begin{frame}
    \frametitle{Technique 1: Set-of-Mark (SoM) Prompting}
    \begin{block}{Definition}
        Overlaying numbered tags or bounding boxes on an image \textit{before} feeding it to the model to improve grounding.
    \end{block}
    \begin{itemize}
        \item \textbf{Problem:} Models struggle to say ``the third cup from the left.''
        \item \textbf{Solution:} You (or a script) draw a red box labeled [1] on the cup.
        \item \textbf{Prompt:} ``Describe the object in box [1].''
        \item \textbf{Result:} Significantly higher accuracy in spatial reasoning.
    \end{itemize}
\end{frame}

\begin{frame}[fragile]
    \frametitle{SoM mini-template (with a tagged image)}
\begin{verbatim}
Context: The image has numbered boxes [1], [2], [3].
Task:
1) For each box, describe what it contains.
2) Answer: Which box contains the power button?
Rules:
- If you are not sure, say "uncertain" and explain why.
Output: JSON
{ "box_descriptions": {...}, "answer": "[n]", "confidence": 0-1 }
\end{verbatim}
\end{frame}

\begin{frame}
    \frametitle{Technique 2: Chain-of-Visual-Thought}
    \begin{itemize}
        \item \textbf{Standard CoT:} ``Let's think step by step.''
        \item \textbf{Visual CoT:} Force the model to break down the visual analysis first.
        \item \textbf{Prompt Structure:}
        \begin{enumerate}
            \item ``First, list all objects visible in the foreground.''
            \item ``Second, describe the action taking place.''
            \item ``Finally, speculate on the relationship between the people.''
        \end{enumerate}
        \item \textbf{Why?} Prevents the model from jumping to a conclusion without ``looking'' at the details.
    \end{itemize}
\end{frame}

\begin{frame}[fragile]
    \frametitle{Template: Extract structured data from a document screenshot}
\begin{verbatim}
Role: You are a careful data entry assistant.
Input: A screenshot of an invoice.
Task:
1) Extract: vendor, date, invoice_id, total, line_items.
2) Quote the exact text for vendor/date/total.
Rules:
- If any field is unreadable, set it to null.
- Do not guess missing numbers.
Output: JSON only.
\end{verbatim}
\end{frame}

\begin{frame}[fragile]
    \frametitle{Prompt Template Example}
\begin{verbatim}
System role: You are an analyst reviewing lab results.
Context: Patient is 42-year-old with Type 2 diabetes.
Attachments: image_1 (glucose chart), image_2 (medication list)
Tasks:
1. Describe key trends in glucose chart.
2. Identify medication conflicts.
3. Highlight missing data to request.
Output format: JSON with keys
{"summary": string, "risks": [string], "follow_up": [string]}
\end{verbatim}
\end{frame}

\begin{frame}
    \frametitle{Combining Modalities}
    \begin{itemize}
        \item Chain modalities: image analysis \textrightarrow{} text reasoning \textrightarrow{} code generation.
        \item Use tables to anchor structured data interpretation.
        \item Provide alt text and plain-language descriptions for accessibility.
        \item Maintain provenance: track which modality informed each claim.
    \end{itemize}
\end{frame}

\begin{frame}[fragile]
    \frametitle{Template: Audio (via transcript) + action items}
\begin{verbatim}
Input: Transcript of a 20-minute meeting.
Task:
1) Summarize in 5 bullets.
2) List action items as: owner, task, due_date (if stated).
3) List open questions.
Rule: Only use information explicitly in the transcript.
Output: Markdown.
\end{verbatim}
\end{frame}

\section{Verification Strategies}

\begin{frame}
    \frametitle{Checking Visual Claims}
    \begin{itemize}
        \item \textbf{Self-Correction:} Ask the model to ``Double check the image. Are you sure the car is red?''
        \item \textbf{Multi-Model Voting:} Ask GPT-4o, Claude 3.5, and Gemini to describe the same image. If they disagree, flag for human review.
        \item \textbf{OCR Validation:} If the task involves reading text, use a specialized OCR tool (like Tesseract or Amazon Textract) to verify the VLM's reading.
        \item \textbf{Provenance Tracking:} Require the model to output bounding box coordinates for its claims (e.g., ``Found a defect at [200, 400]'').
    \end{itemize}
\end{frame}

\begin{frame}
    \frametitle{A quick verification checklist}
    \begin{enumerate}
        \item \textbf{Repeat the question:} Ask it twice with slightly different wording.
        \item \textbf{Ask for evidence:} "Quote the text you read" or "Describe the region".
        \item \textbf{Use a tool when needed:} OCR tool for small text; calculator for arithmetic.
        \item \textbf{Human-in-the-loop:} If the decision is high-stakes, require review.
    \end{enumerate}
\end{frame}

\begin{frame}
    \frametitle{Accessibility Considerations}
    \begin{itemize}
        \item \textbf{Alt Text Generation:} Use VLMs to generate descriptive alt text for millions of images on the web.
        \item \textbf{Video Captions:} Automated speech-to-text + speaker diarization (identifying who spoke when).
        \item \textbf{Visual Interpreters:} Apps like Be My Eyes use VLMs to describe surroundings for blind users.
        \item \textbf{Ethics:} Avoid making assumptions about gender, race, or emotion unless explicitly clear.
    \end{itemize}
\end{frame}

\begin{frame}
    \frametitle{Prompting for Alt Text (practical)}
    \begin{itemize}
        \item \textbf{Short alt text (1 line):} For screen readers.
        \item \textbf{Long description (2--4 lines):} Add context + important details.
        \item \textbf{Avoid sensitive guesses:} Don't infer age, race, disability, or emotions unless clearly stated.
    \end{itemize}
    \begin{block}{Example prompt}
        "Write (1) a one-line alt text and (2) a 3-line long description.
        Do not guess identity attributes. Focus on objective details."
    \end{block}
\end{frame}

\section{Applications}

\begin{frame}
    \frametitle{Industry Use Cases}
    \begin{itemize}
        \item \textbf{E-Commerce:} ``Generate a product description from this photo of a sneaker.'' (Visual Attribute Extraction).
        \item \textbf{Insurance:} ``Analyze this photo of a car accident. Estimate repair costs and list damaged parts.'' (Damage Assessment).
        \item \textbf{Robotics:} ``Look at this messy table. Plan a path to pick up the apple without hitting the cup.'' (Visual Planning).
        \item \textbf{UI/UX Design:} ``Turn this napkin sketch of a website into HTML/Tailwind code.'' (Sketch-to-Code).
    \end{itemize}
\end{frame}

\begin{frame}
    \frametitle{Research and Analysis}
    \begin{itemize}
        \item \textbf{Scientific Charts:} ``Extract the data points from this scatter plot into a CSV file.''
        \item \textbf{Video Analysis:} ``Watch this 1-hour lecture. Give me timestamps for when the professor talks about 'Transformers'.'' (Video Retrieval).
        \item \textbf{Digital Humanities:} Analyzing historical manuscripts or art styles at scale.
    \end{itemize}
\end{frame}

\section{Wrap-up}

\begin{frame}
    \frametitle{Key Takeaways}
    \begin{itemize}
        \item Multimodal prompts expand capability but require careful design.
        \item Describe inputs precisely and control for hallucinations.
        \item Accessibility and ethics must be embedded from the start.
        \item Evaluate new workflows rigorously before deployment.
    \end{itemize}
    \begin{block}{One last rule}
        If the output matters, treat the model like a \textbf{student}: give instructions, examples, and a grading checklist.
    \end{block}
\end{frame}

\begin{frame}{Questions and Discussion}
    \begin{center}
    {\Huge Thank You!}
    \vspace{1cm}

        \textbf{Dr.\ Mehmet Ali Akyol}\\
        \texttt{mehmetali.akyol@gmail.com}
    \vspace{1cm}

    {\large Questions?}
    \end{center}
\end{frame}

\end{document}
